% !TEX root = IE3_expG1_report.tex
\documentclass[lualatex,a4paper,10pt]{jlreq}
% Shared LuaLaTeX preamble.
\usepackage{luatexja-fontspec}
\setmainfont{TeX Gyre Termes}
\setsansfont{TeX Gyre Heros}
\setmonofont{Latin Modern Mono}
\setmainjfont[BoldFont=HaranoAjiGothic-Medium]{HaranoAjiMincho-Regular}
\setsansjfont[BoldFont=HaranoAjiGothic-Bold]{HaranoAjiGothic-Medium}
\usepackage[top=22mm,bottom=22mm,left=23mm,right=23mm]{geometry}
\usepackage{amsmath,amssymb,booktabs,array,longtable,tabularx}
\usepackage{graphicx,xcolor,tikz,circuitikz}
\usetikzlibrary{arrows.meta,positioning,calc}
\usepackage{enumitem,fancyhdr,listings,needspace}
\usepackage[unicode,colorlinks=true,linkcolor=labblue,urlcolor=labteal]{hyperref}
\usepackage{bookmark}
\definecolor{labblue}{RGB}{30,70,112}
\definecolor{labteal}{RGB}{0,100,105}
\definecolor{labgray}{gray}{0.95}
\ModifyHeading{section}{font={\large\sffamily\bfseries\color{labblue}}}
\ModifyHeading{subsection}{font={\normalsize\sffamily\bfseries\color{labteal}}}
\setlength{\parindent}{1\zw}
\setlength{\parskip}{3pt}
\setlength{\emergencystretch}{2em}
\renewcommand{\arraystretch}{1.35}
\setlist{itemsep=3pt,topsep=4pt,leftmargin=2\zw}
\newcolumntype{Y}{>{\raggedright\arraybackslash}X}
\newcolumntype{C}[1]{>{\centering\arraybackslash}p{#1}}
\pagestyle{fancy}
\fancyhf{}
\fancyhead[L]{\small\color{labblue} IE3 後期・電子工学実験}
\fancyhead[R]{\small テーマ\LabCode}
\fancyfoot[C]{\thepage}
\setlength{\headheight}{16pt}
\DeclareOldFontCommand{\rm}{\normalfont\rmfamily}{\mathrm}
\lstset{basicstyle=\small\ttfamily,columns=fullflexible,keepspaces=true,
 breaklines=true,frame=single,backgroundcolor=\color{labgray},
 numbers=left,numberstyle=\tiny,showstringspaces=false,aboveskip=8pt,belowskip=8pt}
\newcommand{\blank}{\rule{22mm}{0.3pt}}
\newcommand{\answerline}{\par\noindent\rule{\linewidth}{0.25pt}\par}
\newcommand{\writing}[1]{\par\Needspace{\dimexpr#1+5mm\relax}\noindent%
 \fcolorbox{labblue!35}{white}{\parbox[c][#1][t]{\dimexpr\linewidth-2\fboxsep-2\fboxrule\relax}{\vspace{2mm}\noindent\strut\hfill\par}}\par}
\newcommand{\hint}[1]{\par\smallskip\noindent{\sffamily\bfseries\color{labteal}記述の視点：}#1\par}
\newcommand{\note}[1]{\par\smallskip\noindent\fcolorbox{labblue!45}{labblue!5}{%
 \parbox{\dimexpr\linewidth-2\fboxsep-2\fboxrule\relax}{#1}}\par\smallskip}
\newcommand{\eqerror}{\varepsilon=\frac{x_{\mathrm{meas}}-x_{\mathrm{th}}}{x_{\mathrm{th}}}\times100\ \%}
\newcommand{\LabSource}{徳山工業高等専門学校 情報電子工学科，
 『2026年度 電子工学実験（3年後期）テキスト』，テーマ\LabCode，
 \expandafter\path\expandafter{\SourcePdf}。}
\newcommand{\reporttitle}{
 \noindent{\small\sffamily\color{labteal}実験報告書・記入ガイド}\par
 \noindent{\LARGE\sffamily\bfseries \LabCode\quad\LabTitle}\par\smallskip
 \noindent 氏名：\blank\quad 班：\blank\quad 実験日：\blank\par
 \note{目的・原理・手順を確認し、実測値と自分の考察を記入する。
 考察のガイドは、検討する事象・使う測定結果・記述の方針を示す。}
}
\newcommand{\prelabtitle}{
 \noindent{\small\sffamily\color{labteal}事前学習・前提知識プリント}\par
 \noindent{\LARGE\sffamily\bfseries \LabCode\quad\LabTitle}\par\smallskip
 \noindent 氏名：\blank\quad 班：\blank\quad 予習日：\blank\par
 \note{用語、回路の動き、計算、測定表への記入の順に読む。
 例題の値は説明用であり、実験結果ではない。}
}
\newcommand{\tocrow}[3]{\hyperref[#1]{#2}& #3 &\pageref*{#1}\\[2mm]}
\newcommand{\documentmap}[1]{
 \section*{目次と概要}
 \noindent 青色の項目をクリックすると該当箇所へ移動する。\par
 {\small\begin{longtable}{@{}p{0.29\linewidth}p{0.61\linewidth}r@{}}
 \toprule {\color{labblue}\bfseries 項目} & {\color{labblue}\bfseries 読むと分かること} & 頁\\\midrule
 \endhead #1\bottomrule
 \end{longtable}}
 \clearpage
}
\newenvironment{terms}{\par\smallskip\begin{longtable}{@{}p{0.23\linewidth}p{0.73\linewidth}@{}}
 \toprule {\color{labteal}\bfseries 用語}&{\color{labteal}\bfseries 意味・回路での役割}\\\midrule\endhead}
 {\bottomrule\end{longtable}}
\newcommand{\term}[2]{\textbf{#1}&#2\\[2mm]}
\newcommand{\guidepoint}[4]{
 \Needspace{32mm}\subsection{#1}
 \noindent{\bfseries\color{labteal}考える事象：}#2\par
 \noindent{\bfseries\color{labteal}参照する結果：}#3\par
 \noindent{\bfseries\color{labteal}記述の方針：}#4\par
}
\newcommand{\reportclosing}[1]{
 \clearpage\section{結論のまとめ方}\label{sec:closing}
 \hint{#1}
 \ConclusionGuide
 \begin{enumerate}
 \item 目的に対応する最も重要な実測結果を、測定条件と数値を添えて述べる。
 \item 原理と一致した点・一致しなかった点を示す。原因は確認済みの事実と推測を分ける。
 \item 実施範囲で言えることと、未確認の条件を一文ずつまとめる。
 \end{enumerate}
 \note{書き出しの型：「○○の条件で△△を測定したところ、□□となった。
 これは◇◇と整合する。ただし××は確認しておらず、一般化には追加測定が必要である。」
 空欄は自分の実測値と根拠で埋める。}
 \noindent 結論（実験後に記入）：\writing{30mm}
 \noindent 改善案・次に確認したい条件：\writing{16mm}
 \section*{参考文献}\label{sec:references}
 \begin{enumerate}
 \item \LabSource
 \item 使用したデータシート・書籍・ウェブ資料（著者、題名、該当箇所、URL、参照日）を追加する。
 \end{enumerate}
}
\newcommand{\prelabclosing}{
 \section*{準備・出典}\label{sec:prelabsource}
 \begin{itemize}[nosep]
 \item 資料、筆記具、記録用紙、電卓と、実験条件・機器設定の記録欄。
 \item 確認問題の解答と、分からない用語・回路点への具体的な質問。
 \end{itemize}
 {\small\noindent\LabSource\par}
}
\newcommand{\simpleflow}[1]{
 \begin{center}
 \begin{tikzpicture}[node distance=5mm and 5mm,
 every node/.style={font=\small},box/.style={draw=labblue,fill=labblue!4,rounded corners=1mm,
 align=center,minimum height=10mm,text width=30mm}]
 #1
 \end{tikzpicture}
 \end{center}
}

\newcommand{\LabCode}{G1}
\newcommand{\LabTitle}{マイコンによるステッピング・モータ制御}
\newcommand{\SourcePdf}{IE3_expG1_A4.pdf}
\newcommand{\ConclusionGuide}{励磁列、実測のステップ角、90度回転の回数を対応させてまとめる。
速度と脱調は観察した間隔・負荷の範囲で述べ、プログラムの指令と機械的な結果を別々に確認したことを示す。\par}
\begin{document}
\reporttitle
\documentmap{
\tocrow{sec:auto1}{1 目的}{励磁列によるモータの角度と方向の制御を実験で確かめる。}
\tocrow{sec:auto2}{2 原理}{励磁列によるモータの角度と方向の制御の仕組みと成立条件を整理する。}
\tocrow{sec:auto3}{3 装置・設定}{機器・定数・測定点を確認し、資料の順序で操作する。}
\tocrow{sec:auto4}{4 方法}{機器・定数・測定点を確認し、資料の順序で操作する。}
\tocrow{sec:auto5}{5 結果}{一周の遷移数、90度回転、脱調条件を表や図に記録する。}
\tocrow{sec:auto6}{6 考察：結果を根拠に説明する}{一周の遷移数、90度回転、脱調条件を根拠に、比較と原因の検討を進める。}
\tocrow{sec:closing}{7 結論のまとめ方}{主要な結果、成立範囲、未確認の条件を短くまとめる。}
\tocrow{sec:references}{参考文献}{使用した資料と外部の根拠を示す。}
}

\section{目的}\label{sec:auto1}
TeCからステッピング・モータの励磁を制御し、1相励磁と1--2相励磁の回転方向、ステップ角、速度の違いを理解する。90度回転と逆転を行うプログラムを作成する。
\section{原理}\label{sec:auto2}
コイルの励磁位置を順に移すと回転子がステップ単位で移動する。資料の配線は
$D_0=A,D_2=B,D_1=\overline A,D_3=\overline B$である。
\par\noindent\begin{tabularx}{\linewidth}{lY}
\toprule 方式 & 出力コード（16進、順方向の例）\\\midrule
1相 & 01, 04, 02, 08\\
1--2相 & 01, 05, 04, 06, 02, 0A, 08, 09\\\bottomrule
\end{tabularx}
逆方向は励磁順序を逆にする。一周に要するステップ遷移数をNとすると$\theta=360^\circ/N$、90度には$N/4$遷移が必要となる。
1--2相では同じモータの角度刻みがおよそ半分となる。
\section{装置・設定}\label{sec:auto3}
TeC、モータ・LEDボード、接続ケーブル、BYOD用PC、資料指定のアセンブラを使用する。
\texttt{OUT 07h}で$D_7$〜$D_0$を出力し、\texttt{0Ch}の上位ビットを1として指定方向に設定する。
モータ型番、励磁配線、一周の実測遷移数、使用ソフトを記録する。
\writing{18mm}
\section{方法}\label{sec:auto4}
\begin{enumerate}
\item 電源OFFでケーブルを接続し、資料指定の電源LED・配線を確認する。通電中の抜差し、軸を手で回す操作を避ける。
\item サンプルを読み、1相励磁を約0.5秒間隔で動作させる。励磁順序と回転方向を確認する。
\item 回転を観察して一周に必要な遷移数を求める。遅延を減らし、脱調の有無を調べる。
\item 1--2相励磁を作成し、同じ確認を行う。90度回転と90度逆転を両方式で作成する。
\item 資料指定の各リストファイルを添付し、待ち時間、繰返し回数、終了時の励磁状態を説明する。
\end{enumerate}
\clearpage
\section{結果}\label{sec:auto5}
表1：励磁方式と回転。
\par\noindent\begin{tabularx}{\linewidth}{YYYYY}
\toprule 方式 & 一周遷移数N & ステップ角/度 & 90度遷移数 & 回転方向\\\midrule
1相&&&&\\1--2相&&&&\\\bottomrule
\end{tabularx}
表2：パルス間隔と動作。脱調は観察した条件を記す。
\par\noindent\begin{tabularx}{\linewidth}{YYYY}
\toprule 方式 & 間隔/ms & 動作・速度 & 脱調・備考\\\midrule
&&&\\&&&\\&&&\\&&&\\\bottomrule
\end{tabularx}
図1：励磁状態の時系列とプログラムの流れ。
\writing{35mm}
プログラムの添付：1相90度、1相逆90度、1--2相90度、1--2相逆90度のソースとリスト。
\writing{20mm}
\clearpage
\clearpage\section{考察：結果を根拠に説明する}\label{sec:auto6}
\guidepoint{励磁方式とステップ角}
{1相と1--2相で一周の遷移数と角度刻みがどう変わるか。}
{一周の実測N、ステップ角の表、励磁コード列と配線対応。}
{$360^\circ/N$を両方式で求め、同じ機械的一周を基準に比較する。原理説明の模式角度と実モータの角度を区別する。}
\guidepoint{90度・逆90度のプログラム}
{指令した遷移数、実際の角度、方向が一致するか。}
{四プログラムのソースと.lst、CALL内の遷移数、回転前後の位置記録。}
{初期状態を別に扱い、CALL回数から遷移数を示す。逆転はコード列の逆順を根拠に説明する。角度を目視だけで読んだ場合は精度を限定する。}
\guidepoint{速度と脱調}
{待ち時間を短くしたとき、指令に追従できなくなる条件はどこか。}
{間隔と動作の表、回転のばらつき、負荷・開始条件。}
{最速の成功条件と最初の失敗条件を示す。急な起動、慣性、負荷を候補にし、未測定のトルクを数値で断定しない。}
\guidepoint{停止後の状態}
{励磁を維持した停止と、解除した停止で保持に違いがあるか。}
{終了時の出力コード、ソース、停止時の観察。}
{どの状態で終了する設計かを述べる。比較を実施していなければ保持力の違いは原理上の予想として示す。}
\subsection{資料の課題・追加検討}
\begin{enumerate}
\item 1相と1--2相の角度、励磁数、動作の違いを説明する。
\item 90度の遷移数を計算し、初期励磁を一回の回転として数えていないか確認する。サブルーチン一回で進む遷移数も示す。
\item 間隔を短くしたときに脱調する理由を、慣性、負荷、加速の観点から述べる。
\item 回転を止めた後の励磁状態と保持の関係を検討する。
\end{enumerate}
\writing{30mm}
\reportclosing{プログラムと実際の角度・方向の対応をまとめる。}
\end{document}
